\section{Classe User.java}
    \subsection{Ruolo}

    \subsection{Codice originale}
    \inputminted[linenos, frame=lines]{java}{../old/User.java}

    \subsection{Analisi di Clean Code}
        \subsubsection{Nomi}      
        I nomi risultano significativi sia dal punto di vista di attributi e variabili sia dal lato delle funzioni.
        L'unico attributo mal nominato risulta essere \texttt{idToken} in quanto rappresenta una ambiguità e non è chiaro se si voglia rappresentare lo UserID o il TokenID del database di Firebase;
        nel primo caso andrebbe rinominato l'attributo, nel secondo andrebbe rimosso in quanto non pertinente con la classe User.
        Un errore più lieve si trova nel nome degli attributi \texttt{name} e \texttt{surname}, è comune utilizzare termini come \texttt{firstName} e \texttt{lastName} 
        per evitare ambiguità in contesti di nomi con multipli termini.
        
        La classe \texttt{User} rispecchia il concetto di utente, seppur con qualche ambiguità dovuta ancora una volta ad all'accoppiamento di diverse responsabilità,
        la classe dovrebbe rappresentare solo l'utente e le sue proprietà senza preoccuparsi di come scambiare i dati tra finestre di Android.

        Altro problema di naming si trova invece nel metodo \texttt{describeContents()}, che risulta un nome esplicativo ma confusionario nell'ambito di utilizzo,
        in questo caso però risulta corretto in quanto parte della implementazione standard dell'interfaccia \texttt{Parcelable}.

        \subsubsection{Responsabilità}
        Come detto la classe viola il \texttt{SRP (Single Responsibility Principle)} in quanto dovrebbe limitarsi a rappresentare un utente e le sue proprietà,
        invece si occupa anche di gestire lo scambio di questi dati tra le differenti attività di Android.

        Andrebbe divisa in due classi distinte, tra cui:
        \begin{itemize}
            \item una classe \textit{User} che rappresenti l'utente e le sue proprietà (nome, email, data di nascita, ecc);
            \item una classe \textit{UserSession} che si occupi dello scambio di informazioni tra le varie finestre e attività di android.
        \end{itemize}
      
        \subsubsection{Formattazione}
        La formattazione risulta corretta, vengono rispettati i limiti consigliati di lunghezza orizzonatale e grandezza delle funzioni; la classe non presenta annidamenti eccessivi, dovuti alla
        mancanza di funzioni molto complesse.

        Il livello di astrazione viene rispettato e l'ordine delle funzioni segue una logica precisa, con setter e getter dello stesso attributo raggruppati insieme; le funzioni più complesse sono spostate
        verso il fondo della classe.

        \subsubsection{Hardcoding}
        La classe presenta un singolo problema di hardcoding, ovvero l'uso della stringa \texttt{"unknown"} per rappresentare un valore di \texttt{surname} non definito.
        Una soluzione migliore sarebbe utilizzare invece il valore \texttt{null}, che rappresenta meglio l'assenza di un valore, se lo si preferisce si potrebbe anche utilizzare una costante statica.

        L'attributo \texttt{age} presenta un problema simile, con il valore iniziale settato a \texttt{0} a rappresentare un'età sconosciuta, questo però comporta ambiguità e problemi di logica.
        Sarebbe consigliato usare un altro valore, come \texttt{-1}, oppure cambiare tipo ad un \texttt{Integer} così da poter assegnare il valore \texttt{null}.

        \subsubsection{Commenti}
        Non sono presenti commenti nel codice della classe, in questo caso si tratta di un fattore positivo in quanto tutto il codice e i nomi delle variabili risultano autoesplicativi.
        Le funzioni sono inoltre brevi, semplice e chiare, non necessitano quindi di commenti aggiuntivi.
        
        \subsubsection{Errori logici}
        La variabile \texttt{age} è adatta ed esplicativa di cosa contiene, ma generalmente è preferibile utilizzare la data di nascita in rispetto all'età, 
        per evitare il costante aggiornamento nella base dei dati.

        Gli errori più gravi si trovano invece nella definizione di setter e getter, in particolare:
        \begin{itemize}
            \item il setter \texttt{setSurname()} non accetta un parametro di input, rendendolo inutile in quanto non è possibile modificare il \texttt{surname} con questo metodo;
            \item il getter \texttt{setAge()} ha lo stesso problema del metodo precedente, non accettando parametri in ingresso non risulta possibile modificare l'età dell'utente.
        \end{itemize}
               
    \subsection{Code Refactoring}
        \inputminted[linenos, frame=lines]{java}{../new/User.java}

    \subsection{Modifiche effettuate}
    Partendo dai nomi, nel code refactoring sono state applicate le seguenti modifiche:
    \begin{itemize}
        \item l'attributo \texttt{name} è stato rinominato in \texttt{firstName} per maggiore chiarezza e standardizzazione;
        \item l'attributo \texttt{surname} è stato rinominato in \texttt{lastName} per maggiore chiarezza e standardizzazione;
        \item l'attributo \texttt{age} è stato rinominato in \texttt{birthDate} e il suo tipo cambiato da \texttt{int} a \texttt{Date}, per rappersentare la data di nascita invece che l'età,
        questo diminuisce la quantità di aggiornamenti necessari e permette di impostare un valore nullo in caso di data mancante;         
    \end{itemize}

    La responsabilità della classe è rimasta invariata, rappresenta ancora un utente e le sue proprietà seppur contenendo anche la gestione delle activity legate ad Android.

    Gli errori in setter e getter sono stati corretti, la classe offre ora la possibilità di modificare gli attributi \texttt{lastName} e \texttt{birthDate} tramite i rispettivi metodi \texttt{set}.

    Il metodo \texttt{toString()} è stato modificato aggiungendo controlli per la data di nascita e per il cognome, in modo da non stampare \texttt{null} ma delle stringhe adatte, 
    andando così a gestire il caso in cui questi siano mancanti o sconosciuti.
       
